\section{Introduction}

Large language models are trained on massive human language corpora, which induce a characteristic geometry in learned representations. Across layers and training stages, activations, gradients, and model parameters exhibit pronounced anisotropy: a small number of directions concentrate disproportionate energy (strong spectral spikes), while the remaining directions form a broad long tail \citep{ethayarajh2019contextual, mu2018all, timkey2021rogue, kovaleva2021bertbusters, machina2024anisotropy}.

While linguistically natural, this structure becomes a computational liability during large-scale training. 
Anisotropy inflates elementwise dynamic range, amplifies extreme activation outliers, and distorts adaptive second-moment statistics 
\citep{kingma2014adam, loshchilov2017adamw, spectra2026}. 
In low-bit regimes, dominant spike directions control quantization scale, compressing long-tail variation into narrow numerical bins 
\citep{xiao2023smoothquant, frantar2022gptq, lin2024awq, li2024svdquant, ashkboos2024quarot}. 
To mitigate these effects, prior approaches rely on explicit spectral control, often via singular value decomposition (SVD) or related orthogonalization procedures, as in Metis 
\citep{cao2025metis}, Muon \citep{jordan2024muon}, and Spectra \citep{spectra2026}. 
Although effective, such techniques are computationally expensive, hardware-unfriendly, and difficult to integrate into large-scale training loops. 
In this sense, anisotropy constitutes a structural curse for efficient model training.

In this work, we examine the internal structure of anisotropy and uncover a sharp asymmetry: anisotropy in activations is dominated by a coherent mean bias. This phenomenon is not accidental. At every representational level---tokens, phrases, clauses, and full sentences---language exhibits shared contextual signal within a given context. Contextual aggregation mechanisms such as self-attention, together with nonlinear transformations, integrate these shared components across tokens. This aggregation naturally induces coherent mean directions in the representation space. As deeper layers capture progressively higher-level semantic regularities, such coherent components reappear and accumulate, producing layer-wise mean drift that propagates through residual connections.

Importantly, due to high dimensionality, even moderate coherent bias produces a consistent directional drift across thousands of coordinates. This rank-one component dominates elementwise extreme values and inflates adaptive second-moment statistics, despite accounting for a small fraction of total variance, making it the primary driver of outlier formation and optimizer instability. Crucially, this structure is not merely harmful---it is also exploitable. Because the dominant instability arises from a rank-one mean component, it can be isolated and removed with a single mean-subtraction operation, which is hardware-friendly and dramatically cheaper than repeated layer-wise spectral decompositions. The same structure that produces instability also exposes a low-dimensional handle for mitigation: the blessing of bias.

We validate this bias-centric interpretation in two downstream settings. 
In low-bit quantization, mean removal substantially reduces dynamic-range inflation and improves long-tail resolution 
\citep{xiao2023smoothquant, wei2023outlier, ashkboos2024quarot, cao2025metis}. 
In adaptive optimization, it contracts extreme second-moment tails, leading to more balanced effective learning-rate allocation 
\citep{kingma2014adam, spectra2026}. 
These results indicate that addressing anisotropy at its structural origin can replace or complement heavy spectral interventions with minimal computational overhead.

\paragraph{Contributions.}
\begin{enumerate}
    \item We identify coherent activation mean bias as a primary structural source of anisotropy in LLM training.
    \item We establish a forward/backward asymmetry: activations accumulate coherent mean drift, while gradient outputs remain mean-neutral.
    \item We quantify that mean-dominated components account for the majority of extreme activation magnitudes and second-moment inflation.
    \item We propose source-level mean removal with optional spike whitening as efficient spectral conditioning mechanisms.
    \item We demonstrate consistent improvements in low-bit quantization and adaptive optimization.
\end{enumerate}
